\subsection{Exact unit norm indices and the level-three split class group}

\begin{statement}{A.175}{Proved}
{The cyclic third step has unit norm index $3^5$ and $S$-unit norm index $3^{10}$}
Put $d=649$, retain the compatible positive roots
$\rho_n^{3^n}=3d^{24}$, and set
\[
 K_n=\Q(\rho_n),\qquad E_n=K_n(\zeta),\qquad
 \zeta=\zeta_3,\qquad G=\operatorname{Gal}(E_3/E_2).
\tag{A.175.1}
\]
Thus $G$ is cyclic of order three.  Let $S$ consist of the places above
$2,3,11,59$ and all infinite places.  Write $U_n=\mathcal O_{E_n}^{\times}$
and $U_{n,S}=\mathcal O_{E_n,S}^{\times}$.  The following are indices of
actual unit groups against actual field norms:
\[
 \begin{split}
 [U_2:U_2\cap N_{E_3/E_2}E_3^\times]&=3^5,\\
 [U_{2,S}:U_{2,S}\cap N_{E_3/E_2}E_3^\times]&=3^{10}.
 \end{split}
\tag{A.175.2}
\]
Moreover, for the action of $G$ on actual ideal classes,
\[
 \left|\bigl(\operatorname{Cl}(E_3)\{3\}\bigr)^G\right|=81,
 \qquad \operatorname{Cl}_S(E_3)\{3\}=0.
\tag{A.175.3}
\]
Here $\operatorname{Cl}_S$ is the quotient of the ordinary class group
by the classes of the finite primes in $S$.  All unit and $S$-unit
elements used to prove these assertions have literal field coordinates
in A.159, A.156, and the coefficient supplement below.  No computed
class-group presentation is used.
\end{statement}

\begin{proof}
\emph{Fields, involutions, and the nine global unit directions.}
Use the actual presentation
\[
 \begin{gathered}
 t^9=3d^6,\qquad \alpha=t^3/d^2,\qquad \rho_2=d^2t,\\
 w=\frac{t^5}{d^3(\zeta-1)},\qquad w^{18}=-3d^6,\\
 \zeta=\frac{w^9/d^3-1}{2},\quad t=-\zeta w^2,\quad
 \alpha=-w^6/d^2.
 \end{gathered}
\tag{A.175.4}
\]
On $E_2$ the involution $\tau$ fixes $t$ and inverts $\zeta$;
it sends $w$ to $-\zeta w$.  Write $\sigma$ for the generator of
$\operatorname{Gal}(E_2/F)$, where $F=K_1(\zeta)$, given by
$\sigma(t)=\zeta t$.  This notation is separate from the group $G$
in A.175.1.  The identities in A.175.4 show that
$\sigma(w)=\zeta^2w$.

Let $u_1,u_2,u_3,u_4$ be the four literal global units of A.159,
with $u_1=U=\alpha^2-2$.  That result proves that they generate the
pro-three completion of $\mathcal O_{K_2}^{\times}$.
For the element $\epsilon$ of A.170, the following formulas fix its
embedding:
\[
 \pi=\frac{\alpha^2}{\zeta-1},\qquad
 v=\frac{1+\pi}{1+\pi^2},\qquad
 \epsilon=v/\tau(v),\qquad v\tau(v)=U.
\tag{A.175.5}
\]
Define $r_2,r_3,r_4$ by the 54 rational coefficients in the
supplement after this proof.  Exact multiplication in
$\Q[w]/(w^{18}+3d^6)$ gives
\[
 r_i^3\sigma^2(u_i)=\sigma(u_i),\qquad
 r_i\tau(r_i)=1\qquad(i=2,3,4).
\tag{A.175.6}
\]
These identities prove that every $r_i$ and its inverse are integral:
their cubes are units and integrality is transitive.  Thus they are
actual global units of norm one for $E_2/K_2$.

Here is a small independent saturation certificate for the five
anti-invariant units $(\zeta,\epsilon,r_2,r_3,r_4)$.
In each row the four entries $(p,t,\zeta,w)$ give a degree-one residue
map of the exact field A.175.4.  The following five entries are the
cubic residue logarithms with respect to the displayed $\zeta$:
\[
 \begin{array}{rrrr|rrrrr}
 p&t&\zeta&w&\zeta&\epsilon&r_2&r_3&r_4\\ \hline
 61&27&47&36&2&2&0&2&0\\
 61&46&47&45&2&2&0&0&1\\
 67&41&37&33&1&0&0&2&1\\
 103&40&46&51&1&0&2&1&1\\
 151&119&32&150&2&2&2&1&2
 \end{array}.
\tag{A.175.7}
\]
The right-hand matrix is invertible over $\F_3$.
Dirichlet's theorem gives rank eight for $U_2$, rank four for its
$\tau$-fixed part, and rank four for its norm-one part.  The
three-primary roots of unity in $E_2$ are precisely $\mu_3$.
Indeed, $\mu_9\subset E_2$ would make $E_2$ contain the splitting
field of $X^9-3d^6$, whose degree is 54: over $\Q(\mu_9)$ its
radical has degree nine.  To check the latter assertion, a cube
root of $3d^6$ in that abelian field would, after multiplication
by a power of $\zeta$, give its positive real cube root $d^2\alpha$.
Its degree-three field is not Galois: the valuation at three proves
irreducibility, while its other two conjugates are nonreal.
It cannot be a subfield of an abelian extension.  The Kummer
degree over $\Q(\mu_9)$ is consequently nine.  This contradicts
$[E_2:\Q]=18$.
The anti-invariant unit group modulo cubes therefore has dimension
five.  A.175.7 proves that the five displayed units form a basis of
it and hence generate its pro-three completion.  The idempotents
$(1\pm\tau)/2$ give the exact primary decomposition.  Consequently
\[
 u_1,u_2,u_3,u_4,\zeta,\epsilon,r_2,r_3,r_4
\tag{A.175.8}
\]
generate the entire pro-three completion of $U_2$.
The finite norm characters below are killed by cubes, so this also
determines their image on the actual unit group.

\emph{The exact local fields and the six tame unit rows.}
At three, $E_2$ has one prime and $v(\rho_2)=2$; adjoining its cube
root is totally ramified of degree three.  At each of eleven and
fifty-nine there are three primes of $E_2$, all with residue field
$\F_{q^2}$.  Each is totally ramified in $E_3/E_2$, since
\[
 v(t)=2,\qquad v(q)=3,\qquad v(\rho_2)=8.
\tag{A.175.9}
\]
The root presentation gives unramified extensions at every other
finite prime: the radicand is a unit there and three is invertible.
All infinite places are complex.  These are therefore exactly the
seven ramified primes in the cyclic extension.

For $q=11,59$, realize the residue field as
$\F_q[z]/(z^2+z+1)$ and send $\zeta$ to $z$.
The three primes are indexed by
$\alpha=a_qz^j$, $j=0,1,2$, where $a_{11}=9$ and $a_{59}=49$.
For a unit $x$, let $\ell_{q,j}(x)$ be defined by
$\bar x^{(q^2-1)/3}=z^{\ell_{q,j}(x)}$.
The tame norm criterion and A.175.9 identify the kernel of this
character with the local norm units.  Its possible nonzero scalar
change upon identifying local Artin groups with $G$ has no effect
on any rank computation below.

The literal residues of A.175.8 are
\[
 \begin{array}{rr|rrrrrrrrr}
 q&j&u_1&u_2&u_3&u_4&\zeta&\epsilon&r_2&r_3&r_4\\ \hline
 11&0&0&0&0&0&1&2&1&2&0\\
 11&1&2&2&1&2&1&0&0&2&1\\
 11&2&1&1&2&1&1&0&0&2&1\\
 59&0&0&0&0&0&2&2&2&1&0\\
 59&1&1&1&2&1&2&1&2&2&0\\
 59&2&2&2&1&2&2&1&2&2&0
 \end{array}.
\tag{A.175.10}
\]
Its rank is five.  The first four columns have rank one and the
last five have rank four.

For completeness the reduction maps used for this table keep the
actual rational coefficients.  If $u=\sum_{i=0}^8c_it^i$ is one of
the first four units, then its residue is
$\sum_{k=0}^2c_{3k}d^{2k}\alpha^k$.  The other two components have
strictly positive valuations, at least two and one respectively
(they are zero for $u_1$).  If $r=\sum_{i=0}^{17}c_iw^i$, put
\[
 h=\frac{\alpha^2(\zeta^2-\zeta)}3,\qquad w^3=dh.
\tag{A.175.11}
\]
Its residue is $\sum_{k=0}^5c_{3k}d^kh^k$; the other components
have valuations at least one and two.  Every denominator left in
these residue formulas is prime to $q$.  The coefficient certificate
checks these bounds term by term, as well as the resulting nonzero
residues.  A.175.5 gives the residues of $\epsilon$ directly.

The local norm characters of a global unit vanish at all other
places except possibly three.  At an unramified finite extension
all units are norms.  Global reciprocity makes the character at
three the negative sum of the six characters above, after their
actual local Artin identifications with $G$.
The Hasse norm theorem for this cyclic extension
\cite[Chapter VIII, Theorem 3.1, printed p.~234]{MilneCFT}
therefore identifies
the kernel of the six-row map with $U_2\cap N E_3^\times$.
The complete generating set A.175.8 and rank five in A.175.10 prove
the first equality of A.175.2.

\emph{Four additional valuation rows and a full $S$-unit image.}
At two, $K_2$ has the three unramified factors of degrees $1,2,6$
described in A.165.  Adjoining $\zeta$ makes the degree-one factor
one degree-two completion, and splits each of the other two factors
into two completions.  The unique first completion splits completely
in $E_3/E_2$.  Each of the other four has an unramified extension of
degree three.  These statements also follow directly from the
successive cubic-root factors: over the old rational branch the
three roots become rational after adjoining $\mu_3$, while the
other two old branches acquire an unramified cubic factor.
More explicitly their residues have orders three and nine in
$\F_4^\times$ and $\F_{64}^\times$, respectively; a next cube root
has order nine and twenty-seven.  The exact orders
$\operatorname{ord}_9(2)=6$ and $\operatorname{ord}_{27}(2)=18$
give the indicated cubic residue-degree increases.  Principal units
have unique cube roots, so these residue calculations identify the
actual unramified local extensions.
Thus exactly four primes above two have nontrivial local norm
quotient, and its character is valuation modulo three, up to the
nonzero Frobenius scalar.

Let $\beta_1=1-\alpha$ and let $\beta_2$ be the literal element of
A.156, with relative norm $\beta_1$ and divisor the distinguished
degree-one prime of $K_2$ above two.  Adjoin to A.175.8 the five
actual $S$-units
\[
 \beta_2,\quad s_2=\beta_1/\beta_2,\quad
 s_3=2/(\alpha-1),\quad
 a_{20}=\frac{\sigma(\beta_2)}{\sigma^2(\beta_2)},\quad
 a_{26}=\frac{\zeta\alpha-1}{\zeta^2\alpha-1}.
\tag{A.175.12}
\]
All five are units above eleven and fifty-nine.  The six cubic
residue logs of $\beta_2$ are, in the row order of A.175.10,
\[
                   (0,2,1,0,2,1)^t.
\tag{A.175.13}
\]
They are obtained from the literal A.156 coefficients by the same
three-component formula as for $u_i$; the discarded components
again have valuations at least two and one.
This column raises the rank of A.175.10 to six.  The other four
additional columns are respectively its negative, zero, zero, and
zero.  For $a_{20}$, residue triviality of $\sigma$ at a ramified
prime gives residue one; the remaining entries follow by the
displayed rational formulas.

The exact divisors supply the independent two-adic block.
Label the four growing primes $P_{2,+},P_{2,-},P_{6,+},P_{6,-}$.
The first pair lies over the old degree-two prime of $K_2$, and
the second over its old degree-six prime.  The valuation block on
$(s_2,s_3,a_{20},a_{26})$ is
\[
 \begin{array}{c|rrrr}
 &s_2&s_3&a_{20}&a_{26}\\ \hline
 P_{2,+}&1&0&1&0\\
 P_{2,-}&1&0&-1&0\\
 P_{6,+}&0&1&0&1\\
 P_{6,-}&0&1&0&-1
 \end{array}.
\tag{A.175.14}
\]
Here the first two columns are the actual divisors in A.165.
To identify the third, the old two-prime of $F$ supporting
$\beta_2$ splits into three primes in $E_2/F$; $\sigma$ permutes
them cyclically and $\tau$ fixes the distinguished one and swaps
the other two.  The quotient defining $a_{20}$ has valuations
$1,-1$ on the latter pair.  The numerator and denominator of
$a_{26}$ generate the other two norm-four primes of $F$.
Those primes acquire unramified degree-three extensions in $E_2$,
giving precisely its two degree-six primes with valuations $1,-1$.
This proves the full literal block A.175.14.
The nine global units and $\beta_2$ have valuation zero at these
four primes.  The block has rank four modulo three, and
A.175.10 together with A.175.13 has rank six.  The resulting
fourteen-column, ten-row local norm matrix therefore has rank ten.

There are exactly eleven nonsplit places of $E_2$ in $S$: four
above two, one above three, three above eleven, and three above
fifty-nine.  Each local quotient is $G\simeq\Z/3$.
For an $S$-unit every local character outside $S$ vanishes, and
global reciprocity places the eleven characters in the kernel of
the sum map $G^{11}\to G$, of order $3^{10}$.
The ten independent rows just proved give the reverse index bound.
The cyclic Hasse norm theorem identifies the kernel with
$U_{2,S}\cap N E_3^\times$.  This proves the second equality
of A.175.2.  This argument does not require an unproved claim that
the fourteen exhibited elements generate the entire $S$-unit group.

\emph{The precise ambiguous-class formula, including the $S$ factors.}
We give the cohomological comparison used here.  For any cyclic
extension $L/K$ with group $G$, use the groups of ideals and
principal ideals away from $S$, denoted $I_{L,S}$ and $P_{L,S}$.
The exact sequences
\[
 \begin{gathered}
 1\longrightarrow U_{L,S}\longrightarrow L^\times
       \longrightarrow P_{L,S}\longrightarrow1,\\
 0\longrightarrow P_{L,S}\longrightarrow I_{L,S}
       \longrightarrow\operatorname{Cl}_S(L)\longrightarrow0
 \end{gathered}
\tag{A.175.15}
\]
retain the actual divisor maps.  Hilbert 90 gives
$P_{L,S}^G/P_{K,S}=H^1(G,U_{L,S})$.
The permutation ideal module has $H^1(G,I_{L,S})=0$, and cyclic
two-periodicity gives
\[
 H^1(G,P_{L,S})
       =\frac{U_{K,S}\cap N L^\times}{N U_{L,S}}.
\tag{A.175.16}
\]
Extension of ideals has index
$[I_{L,S}^G:jI_{K,S}]=\prod_{v\notin S}e_v$.
Taking indices in A.175.15--16, including the possibly nontrivial
capitulation encoded by $H^1(G,U_{L,S})$, yields
\[
 |\operatorname{Cl}_S(L)^G|
  =h_{K,S}\left(\prod_{v\notin S}e_v\right)
       \frac{h_G(U_{L,S})}
       {[U_{K,S}:U_{K,S}\cap N L^\times]}.
\tag{A.175.17}
\]
Here $h_G$ is the cyclic Herbrand quotient
$|\widehat H^0|/|H^1|$.
For the totally imaginary cyclic extension in this proof, the
rational Dirichlet module is
$\Q[G]^{r_2(K)}-\Q$; its quotient is $1/3$.
Adding an $S$-prime adds the permutation module $\Z[G/D_v]$
up to finite index.  Its Herbrand quotient is $|D_v|$, by Shapiro's
lemma and the trivial integral module.  Finite modules have quotient
one.  Consequently
\[
 \begin{gathered}
 h_G(U_{L,S})=\frac13\prod_{v\in S_{\rm fin}}|D_v|,
 \\
 |\operatorname{Cl}_S(L)^G|
 =\frac{h_{K,S}
     \prod_{v\notin S}e_v\prod_{v\in S_{\rm fin}}|D_v|}
        {3[U_{K,S}:U_{K,S}\cap N L^\times]}.
 \end{gathered}
\tag{A.175.18}
\]
Complex places have decomposition order one and add no factor.
The formula also applies with no finite primes in $S$.

Apply it first to ordinary classes of $E_3/E_2$.
A.173 gives $|\operatorname{Cl}(E_2)\{3\}|=27$; the seven
ramified primes and the first norm index give
\[
 \left|\bigl(\operatorname{Cl}(E_3)\{3\}\bigr)^G\right|
                   =\frac{3^3\,3^7}{3\,3^5}=3^4.
\tag{A.175.19}
\]
For the stated $S$, all ramification is inside $S$,
the eleven nonsplit places give $\prod|D_v|=3^{11}$,
and A.173 proves $\operatorname{Cl}_S(E_2)\{3\}=0$.
Thus
\[
 \left|\bigl(\operatorname{Cl}_S(E_3)\{3\}\bigr)^G\right|
                       =\frac{3^{11}}{3\,3^{10}}=1.
\tag{A.175.20}
\]
A nonzero finite three-group with an order-three automorphism has
a nonzero invariant: on its nonzero subgroup killed by three,
$\sigma-1$ is nilpotent over $\F_3$ and has nonzero kernel.
Equation A.175.20 therefore proves the asserted vanishing of the
entire three-primary $S$-class group.
\end{proof}

The exact coefficient identities and saturation rows are independently
checked by \path{code/two_inert_level2_anti_cube_root_certificate.py}.
The entire unit and $S$-unit norm matrix, including every discarded
component valuation and the four-prime divisor block, is checked by
\path{code/two_inert_level2_unit_norm_certificate.py}.
Their executed outputs and literal data are retained.
The coefficient supplement is generated and reparsed by
\path{code/write_two_inert_level2_anti_coefficients.py}.
These programs verify exact field and residue identities; the class-group
conclusions follow from the proved index and cohomology maps above.

\subsection*{Literal anti-unit roots used in A.175}
Put $d=649$, $B=10^{40}$ and $w^{18}=-3d^6$ as in A.175.4.
The following integers specify the three elements without root-choice ambiguity:
\[r_i=\frac1{2d^5}\sum_{j=0}^{17}n_{i,j}w^j,\qquad
n_{i,j}=\epsilon_{i,j}\sum_{k\geq0}b_{i,j,k}B^k\quad(i=2,3,4).\]
Each row gives two successive forty-digit blocks; unlisted blocks are zero.
The sign is repeated on every row. Exact multiplication proves
$r_i^3\sigma^2(u_i)=\sigma(u_i)$ and $r_i\tau(r_i)=1$.
\begingroup\setlength{\tabcolsep}{3pt}\fontsize{7.6}{9.2}\selectfont\ttfamily
\par\medskip\normalfont\normalsize Root $r_2$\par\smallskip
\fontsize{7.6}{9.2}\selectfont\ttfamily
\begin{longtable}{rrrll}
$j$ & $k$ & $\epsilon$ & $b_{i,j,k}$ & $b_{i,j,k+1}$ \\ \hline\endhead
0 & 0 & + & 0504427827433875490862791429022420320508 & 0000000075228655954968336734523604987140 \\
1 & 0 & + & 8045486861039111729335670540409669690004 & 0000000005955763964715263970167366545772 \\
2 & 0 & + & 3594896215524231153925054198950584602659 & 0000000000232708942764489301228617579549 \\
3 & 0 & - & 1539658390758981677456551273663809623465 & 0000000000000021187781706628821329234152 \\
4 & 0 & - & 6152616032234897364346584541615670221477 & 0000000000004722198612579045338060312672 \\
5 & 0 & - & 7074057935481831349695032332293520011099 & 0000000000001054980215274212567578834848 \\
6 & 0 & - & 7896624664280648514534804345899228651638 & 0000000000000107158406302678830209048600 \\
7 & 0 & - & 3509471957611368633071007781163013033839 & 0000000000000008394979182370902530274712 \\
8 & 0 & - & 1771716503238600724491062442012136417402 & 0000000000000000684755639219057086617059 \\
9 & 0 & + & 5062076492528969604824199954894890674638 & 0000000000000000005820495754004496971537 \\
10 & 0 & + & 6654267875699746937494359513371452946068 & 0000000000000000010905510163799991693572 \\
11 & 0 & + & 3732926054186907778333612523281131292921 & 0000000000000000001435542317450642970343 \\
12 & 0 & + & 2546860282644904407963620743358158158313 & 0000000000000000000161799342778371038816 \\
13 & 0 & + & 7346685397101338899787102553731648508889 & 0000000000000000000018460034618191179177 \\
14 & 0 & + & 0986512130119408722955188109206856839003 & 0000000000000000000000954665355265811571 \\
15 & 0 & - & 6062157563823043120057866332662335060810 & 0000000000000000000000051550020176061325 \\
16 & 0 & - & 5519366087459106395093588718388620755433 & 0000000000000000000000012795296698333404 \\
17 & 0 & - & 6315572181614970161620480587886704728390 & 0000000000000000000000002250883862523281 \\
\end{longtable}
\par\medskip\normalfont\normalsize Root $r_3$\par\smallskip
\fontsize{7.6}{9.2}\selectfont\ttfamily
\begin{longtable}{rrrll}
$j$ & $k$ & $\epsilon$ & $b_{i,j,k}$ & $b_{i,j,k+1}$ \\ \hline\endhead
0 & 0 & + & 2188849738344761989493946294191332221412 & 8978982852997270502960526008784583692347 \\
0 & 2 & + & 2662745438212591433101343496352071987845 & 4023154723028614451130491803276694505150 \\
0 & 4 & + & 9188508944710405627489127252045757259896 & 1842930255318901744952091955745000808837 \\
0 & 6 & + & 0046528190580462785308848976703106267201 & 0000000000000000000000000000000000000000 \\
1 & 0 & + & 6145547559520951201025952098959402518266 & 9881840616865140498651539447346118345924 \\
1 & 2 & + & 5985879610187446401480256995069226523318 & 2231721603994113330661075361204346224896 \\
1 & 4 & + & 5655244740025115761579431493368861805296 & 6105793984764302610707393928268131592972 \\
1 & 6 & + & 0005747108911700409046451435501198967422 & 0000000000000000000000000000000000000000 \\
2 & 0 & + & 0706895583087103211976441866208292132598 & 3581611446549669019575107335460082782752 \\
2 & 2 & + & 0789309457342197285888393204700240071891 & 5681524886832523116123689399109064431490 \\
2 & 4 & + & 8861899775304363228233842183318135951567 & 7163716816255247228884447231986976453353 \\
2 & 6 & + & 0000680631651082057826890700556810213433 & 0000000000000000000000000000000000000000 \\
3 & 0 & + & 3861623734670038820455106162231817284237 & 8686139027126468376903308176490054838864 \\
3 & 2 & + & 4853860523358573424186430682851380119769 & 3341974245316515470839843093806695898724 \\
3 & 4 & + & 0806265494241635981170971394856912331826 & 5056774650315046692900108452138964226371 \\
3 & 6 & + & 0000077811820119225494460485173627075399 & 0000000000000000000000000000000000000000 \\
4 & 0 & + & 0084581162832684355421214307283707677670 & 5030995434481093696577316177262933542728 \\
4 & 2 & + & 4610305763769461975859029429912386030917 & 0719236980248705251806422212644453327370 \\
4 & 4 & + & 0969980797227960288225459603427377872930 & 7079130299869012549224795310813358355930 \\
4 & 6 & + & 0000008616959540829262651405383173615787 & 0000000000000000000000000000000000000000 \\
5 & 0 & + & 8573018104638004905613757733114612147546 & 4521788605190059562816842204354563390256 \\
5 & 2 & + & 7501830054423852749494168657375499319020 & 1345554533404892442336998256226706277028 \\
5 & 4 & + & 8898591660780700423921246595907186719378 & 3582853013490044905566471595102055170199 \\
5 & 6 & + & 0000000925464390083845077690915212661370 & 0000000000000000000000000000000000000000 \\
6 & 0 & + & 3601519143273206931197684300394122223887 & 8759281590715344482079325968479713505904 \\
6 & 2 & + & 3626795845561972559083937262484390166136 & 6241284255172425771698457437186655435428 \\
6 & 4 & + & 9947169755195257845186406567259484253854 & 3765621667079945048209623575401441150185 \\
6 & 6 & + & 0000000096325915910534508434908166808860 & 0000000000000000000000000000000000000000 \\
7 & 0 & + & 3700233840496534114134579324954431498460 & 0904079313852603188139153104663526389709 \\
7 & 2 & + & 5805430056421252577150611014470407911052 & 1931796607420127449762646303208018098640 \\
7 & 4 & + & 4686067543906176911307114337923850674153 & 6640643910337985004388833854229678443778 \\
7 & 6 & + & 0000000009688541721844583287263620666468 & 0000000000000000000000000000000000000000 \\
8 & 0 & + & 1045083098128022804210544700702482161338 & 2595658789508273878864912542501915182600 \\
8 & 2 & + & 6664826199908078433807231714280479632638 & 0342074091053334959709657298370407796560 \\
8 & 4 & + & 9400312011157899880853190168527911929430 & 8178862939514569279991931707572801459190 \\
8 & 6 & + & 0000000000936202722049653882770423905288 & 0000000000000000000000000000000000000000 \\
9 & 0 & + & 7711195312608772483589346566966965354440 & 8745210021320085479321457595734635059468 \\
9 & 2 & + & 3440679443238524259524166734181731418236 & 3758170640655424458714106284081687937159 \\
9 & 4 & + & 4694754320930878441071558756683651977004 & 0748173283312137869380685614994022727764 \\
9 & 6 & + & 0000000000085971491692074293388425488255 & 0000000000000000000000000000000000000000 \\
10 & 0 & + & 9689857434394670776208121005259334713296 & 4598442738279884064724976961034821749336 \\
10 & 2 & + & 5353806932267228585590742816914307438487 & 9816370485276293579663957877390666745454 \\
10 & 4 & + & 2795674501822107259800569070619640799799 & 3368571901793029713595899539880694912327 \\
10 & 6 & + & 0000000000007345667494879554729046280107 & 0000000000000000000000000000000000000000 \\
11 & 0 & + & 9404303606548557907331302980331163881218 & 7836799523319096144694924869759057849367 \\
11 & 2 & + & 2700147517647364449290318486878159571480 & 7779341228175494555630878337048199756330 \\
11 & 4 & + & 7590600320474941790537512467669201772512 & 0415266542875850414392915115545760680763 \\
11 & 6 & + & 0000000000000557033995621131346930084274 & 0000000000000000000000000000000000000000 \\
12 & 0 & + & 1895589437659424714399395072273010488707 & 2514254776692773083618001352668986408495 \\
12 & 2 & + & 5541738360310722390222706628481242501173 & 0460219253697125566359849915449610254032 \\
12 & 4 & + & 7943259834544693651689498374263379798899 & 5465299628375699703887203573385660874232 \\
12 & 6 & + & 0000000000000032484190120288970205780943 & 0000000000000000000000000000000000000000 \\
13 & 0 & + & 5609942125044200441887283256159620258636 & 6969717992024952524816790496607805073130 \\
13 & 2 & + & 6357465397027064575833776190754074338538 & 4033533977249749755034251160161328587370 \\
13 & 4 & + & 1919615289770725756077649558531690752144 & 2857349183666800487282751674837513359661 \\
13 & 6 & + & 0000000000000000375193415083709681604971 & 0000000000000000000000000000000000000000 \\
14 & 0 & - & 2019448543416978089587480296164185690128 & 8397890671881359288646418461295818869264 \\
14 & 2 & - & 5583061431196250055640762411738742999596 & 5403822840304143249544606154994127759645 \\
14 & 4 & - & 1602255645654701582241799276629672670811 & 8641265446207295917765999425369778382799 \\
14 & 6 & - & 0000000000000000303256088925323541821175 & 0000000000000000000000000000000000000000 \\
15 & 0 & - & 9708676180668014453564730359155549514101 & 5299153746820536114768970106232681097693 \\
15 & 2 & - & 1272241636341829127635582176584227775276 & 8093931452026033388152743062430568866731 \\
15 & 4 & - & 9328029296629226083550444839184115955764 & 2941137432533328216140022821098709685597 \\
15 & 6 & - & 0000000000000000069333767818412834204153 & 0000000000000000000000000000000000000000 \\
16 & 0 & - & 0742640522206915977995339402999980832172 & 9380582598565634477823252751356886129100 \\
16 & 2 & - & 3296932767662720964092174121261468388829 & 5369891388933657259044396352250290722046 \\
16 & 4 & - & 6965804204553597021064983981938919982951 & 7779606759504605590079496398924240887164 \\
16 & 6 & - & 0000000000000000011258313360251013256911 & 0000000000000000000000000000000000000000 \\
17 & 0 & - & 0485415596938946759659341767416714954580 & 0482544394524853214590623326423982951899 \\
17 & 2 & - & 8409181077942084066847882528422967052471 & 7086499757041573087964090872452281465622 \\
17 & 4 & - & 3518400075932600034060642971380999069068 & 2897256460905976795629988386529366423863 \\
17 & 6 & - & 0000000000000000001590878361307837272906 & 0000000000000000000000000000000000000000 \\
\end{longtable}
\par\medskip\normalfont\normalsize Root $r_4$\par\smallskip
\fontsize{7.6}{9.2}\selectfont\ttfamily
\begin{longtable}{rrrll}
$j$ & $k$ & $\epsilon$ & $b_{i,j,k}$ & $b_{i,j,k+1}$ \\ \hline\endhead
0 & 0 & + & 8174761203488695439473712343827600186171 & 1770566437160271828688635256011333623211 \\
0 & 2 & + & 5065542037948102497539402121653345944957 & 7785987967080900610030757173747661230138 \\
0 & 4 & + & 0780285448905529193371455579848492020222 & 1875489231217106967459696389631705118022 \\
0 & 6 & + & 7233380282032343173871699742463410422650 & 4189548681416991960048164495181982129452 \\
0 & 8 & + & 2464877815456740930375889608407930386493 & 0289090657855257108201620640828286100556 \\
0 & 10 & + & 8849435437522255020463759596948793479615 & 1015323907742439368377164939171287195117 \\
0 & 12 & + & 7023965410072146991193120657979179563331 & 2582761877355712201880773162532333974904 \\
0 & 14 & + & 2112691238678002769979165603259498594338 & 1937139504439595572213536604858975135456 \\
0 & 16 & + & 3378352865912418984151073926334258790428 & 0192023458125288641867679553046582660603 \\
0 & 18 & + & 4894304463510095456528338263290933461582 & 4949796340418506445225148052250268186469 \\
0 & 20 & + & 1340165635218239326726719152281441674233 & 4057906050990318387621029734281731074345 \\
0 & 22 & + & 4383265880504479304444911987483659534837 & 2905579246214104593783735263959145088626 \\
0 & 24 & + & 7192352634397222265526532715796492064281 & 4742494219611490302296045854574158444805 \\
0 & 26 & + & 4393186237515706814711531925037585336514 & 0028056309975580429176400811978333273505 \\
0 & 28 & + & 2432122168235565472258319133053258806565 & 0913134821788994716078145202963389432935 \\
0 & 30 & + & 1026448596318889324314881892918299022089 & 4584222146177126776704715618607647500835 \\
0 & 32 & + & 9438859526631586729720238128986142551340 & 3821463260774754294177562458734385558480 \\
0 & 34 & + & 8552413688707743559690691350263353569809 & 5754823349987329286623353255148616790432 \\
0 & 36 & + & 5970429953982983921872045345715683569953 & 8550981795572634464467440252608009254876 \\
0 & 38 & + & 4442596105219300091585490859749547907733 & 7235129638113248109320291451311166360960 \\
0 & 40 & + & 0000296872749754757281479972522566294513 & 0000000000000000000000000000000000000000 \\
1 & 0 & - & 4826588981349678069215895358013483721619 & 3441164467655992950982676907924991076052 \\
1 & 2 & - & 7877615840628953712123650657789108481892 & 9022702756571751201010776997865560790486 \\
1 & 4 & - & 2299011167625201009289066276091082427169 & 2199541277543707779980341374217552237885 \\
1 & 6 & - & 3909309680712194690997947233953369806298 & 9316502212391351444473562363131285137689 \\
1 & 8 & - & 2371008378745748760321974498306666262147 & 5601463054022318611914784623490784742608 \\
1 & 10 & - & 8697751829859445430351757283302653840117 & 4865777657255179778690058397514742804147 \\
1 & 12 & - & 5214242859776534286836374658376297650614 & 0158100824262732549710910621388531522991 \\
1 & 14 & - & 0548349929630049817872107183877628064958 & 7096189184031108879739386631323602597757 \\
1 & 16 & - & 4220608081010919628846157806055064698947 & 4325730394129798743274168324248963662024 \\
1 & 18 & - & 0786950380884225403603302684019440553547 & 0790753890268217076598875643296552568161 \\
1 & 20 & - & 0242060312264854520041084233688666839312 & 4786066931631724080435601073504712935620 \\
1 & 22 & - & 6447472276290509986756867012674378202003 & 8826008953100461260629667284953681938424 \\
1 & 24 & - & 2845289364669047835295258619644263280685 & 2537162040800545478893848460885506645749 \\
1 & 26 & - & 1259865599697559697501718955125018989475 & 2769434281422961498283918857104088172087 \\
1 & 28 & - & 2355448793954334468499935654057874094112 & 8364761846029920728366807074420844945666 \\
1 & 30 & - & 8061151042503308959939029024911343441609 & 3294649724519330638279670011215890335920 \\
1 & 32 & - & 7296087034397827706856102281409481685641 & 0988550905800806042674311748640152707377 \\
1 & 34 & - & 2222234040247339962685004252003111121433 & 7011243370399573938735372262487839932956 \\
1 & 36 & - & 4085443978414524566031702946996903502698 & 4459878237507199635420837897956712120285 \\
1 & 38 & - & 3555231977348879265906380081681304456120 & 3814176383334269984363777713524204417696 \\
1 & 40 & - & 0000039694116196358544616151373087854391 & 0000000000000000000000000000000000000000 \\
2 & 0 & - & 0193841673265405478425934736890881799782 & 8410320301622930455960394121582549958284 \\
2 & 2 & - & 4843684058272333579285299036294173241144 & 8877561545690858452137636099270583080037 \\
2 & 4 & - & 8554498939507521482661057981904516886608 & 5125271417548007034766419013161194022885 \\
2 & 6 & - & 7840252335773615220913316629813512403870 & 5684902682342516453890892210902837429629 \\
2 & 8 & - & 5226284332199820048456403033768747173182 & 3248869469905631680063986253129575105815 \\
2 & 10 & - & 1943208987884424536753699409541853204019 & 9767826560622975593547666233486230501064 \\
2 & 12 & - & 0980006921662198749678123032552206764463 & 0531918685914563408987344133645167292241 \\
2 & 14 & - & 4384086364362575586758709883220159283466 & 0056200184503725361556025844525010182768 \\
2 & 16 & - & 2569068310632948160857912258526420904082 & 7747638800530643663787508633392502122383 \\
2 & 18 & - & 6868475094770483889751092745730409000818 & 0355638717883783749658940579114886546371 \\
2 & 20 & - & 9571652225800429406304551786871208878411 & 8481681785845065052952914343447392008451 \\
2 & 22 & - & 9889954817952595015382627787153681331517 & 4814896988119992277593749004269276980117 \\
2 & 24 & - & 6758914762775522532942224275150293790758 & 8190263519081156213763069548672677323388 \\
2 & 26 & - & 1329880921840649817674811617209913905726 & 9217317370915217201889539224543722536781 \\
2 & 28 & - & 8470872362853426799781424178840017038817 & 3329644573972470993887038393376603442679 \\
2 & 30 & - & 7256401701032932335163299657529256981730 & 1916159628351214454714826099289988370860 \\
2 & 32 & - & 3346016419399558325844669174014355807840 & 1750700543380119169513447518216284457696 \\
2 & 34 & - & 0350922518659004421892897157199507573332 & 4472241910117103018094765609164853373732 \\
2 & 36 & - & 6414876502977759272764812945580885714215 & 6568336341150475956938666875485761976949 \\
2 & 38 & - & 9660764327881759036072434496831149322053 & 6975043832113781262558730818837293695041 \\
2 & 40 & - & 0000010976081022350430505591518796872019 & 0000000000000000000000000000000000000000 \\
3 & 0 & - & 2122074388867582952572745382186088050662 & 9049699518141269037045999031485813803614 \\
3 & 2 & - & 9814272954485663078209925485598799089274 & 5010680143041034381069670451625124817911 \\
3 & 4 & - & 6125022207578962756833600213656218190622 & 8054082393019044972178784675341726923686 \\
3 & 6 & - & 5658781166926455068812244564699367354774 & 7439060803814315885131885939528245714074 \\
3 & 8 & - & 0363921093424066797394469431280238489260 & 8053358136123054215731659794128729304029 \\
3 & 10 & - & 5385707394514464046753600703124435581247 & 9856539599867871092815724634725885359806 \\
3 & 12 & - & 0012811479736418633122113690200229198482 & 0764112211155940468759201194489920983482 \\
3 & 14 & - & 5356043943998904627839019745301970156222 & 7851668375074171197127291958696858476434 \\
3 & 16 & - & 6430259931622717632282127975944220203316 & 0125354681180511214791756804014836128104 \\
3 & 18 & - & 0741675678688632683642231995244999898781 & 3312731097510972913834871650292027651382 \\
3 & 20 & - & 2720283248655780578894523564072368151136 & 2311909955721246236938258572483037440332 \\
3 & 22 & - & 2910343464915761666682109451392568863934 & 3179994032315061955989894508684409596911 \\
3 & 24 & - & 9130346537785015105893794168900932256491 & 2578817852489304322038739429571628707716 \\
3 & 26 & - & 5236514848439237078330029604927477107613 & 9307341635337175277720870278228916915425 \\
3 & 28 & - & 8793411160614861824329331315911477218123 & 1677620441438953202623302632044616610049 \\
3 & 30 & - & 9376798097784851060253821044108093302675 & 2526298575465711608057849465883123656069 \\
3 & 32 & - & 9783392191328537951883905937071136061336 & 2282537210935010771126487702695635155505 \\
3 & 34 & - & 3899230005778642259071023959414661075416 & 9508969017441882557860526167465995966105 \\
3 & 36 & - & 1329950549125522548385219740384000340369 & 5357779820637508814825702367832373874171 \\
3 & 38 & - & 3615905320080713744439806197203349173115 & 1730170896579634758517830820276889572804 \\
3 & 40 & - & 0000001597103641791611055082598339895049 & 0000000000000000000000000000000000000000 \\
4 & 0 & - & 7878624686122969429676810908064523291411 & 4278257305439998082183798818472800626351 \\
4 & 2 & - & 6490191418033871743734106373614594225270 & 4596304140392930562334160359775646001854 \\
4 & 4 & - & 8067027343103978154323409701856402357597 & 6179877831700288471120321109200812853922 \\
4 & 6 & - & 2686189647778431856922078454474335118902 & 2591785382370902496992140821387886510129 \\
4 & 8 & - & 9671732984158333133202251315997099284865 & 1909282501217905657503464467956219042305 \\
4 & 10 & - & 4130663219292125879476201222314250999884 & 9136638462176436505563697277939524479827 \\
4 & 12 & - & 0608044491589171696176346336280013652398 & 7020951569648522660631454428951077975045 \\
4 & 14 & - & 3324296952684566151866299816671040970452 & 5705357404044529477225573264205733504341 \\
4 & 16 & - & 6808644223299366881998429761373279835755 & 3707856158651762911679599289342295484564 \\
4 & 18 & - & 9887618204956285154965605828142535887066 & 5890168734357760771514961016668922422870 \\
4 & 20 & - & 1375314414152864445203003840834469406573 & 6342690351141509878515704176159110080498 \\
4 & 22 & - & 6640915979327762862995488439102725467569 & 7865536819182434659779178330680758663884 \\
4 & 24 & - & 2000749846961883973477618225701809083432 & 2375621006858558933891217345814578615110 \\
4 & 26 & - & 8084030963810461372019931974099470747170 & 6573464470095098210334807971757395270915 \\
4 & 28 & - & 8227439431450005542728023089482697165099 & 8400787106148585501879195153945579975300 \\
4 & 30 & - & 5156025303091523986187239611775344578639 & 0589920011016405388416044594779760970910 \\
4 & 32 & - & 4410688548371989375336933421191655601589 & 9310613716055505887550066558380416935101 \\
4 & 34 & - & 1827309859231686211124549450783372221846 & 5017474112008498222714502986130845872286 \\
4 & 36 & - & 1892886962527727198021591035795034879047 & 4881069026943590163286924291001913885931 \\
4 & 38 & - & 7828824935157801283269356518149577633511 & 4735069282195319281445751463113996574370 \\
4 & 40 & - & 0000000170988602543094446787741421214067 & 0000000000000000000000000000000000000000 \\
5 & 0 & - & 7855737256531277905908730179367476866907 & 3159971330346099851785213304246021820244 \\
5 & 2 & - & 0906022500384270833282160494255699768712 & 3558747047198586798703378993636224633364 \\
5 & 4 & - & 7364900067426763715551098889238496782047 & 8753819592587772086192899049460044031371 \\
5 & 6 & - & 8081238987476862892441119904083797118386 & 5897560998950112391591801240532003289946 \\
5 & 8 & - & 8303345023742723016513707694986377112409 & 3634866944831004855697697588516951895422 \\
5 & 10 & - & 4532745764150203091529194716499405018699 & 7046858104029606005634506363946934583601 \\
5 & 12 & - & 1348648560698473234228888341693649459688 & 2990724779135103379834888350628858794060 \\
5 & 14 & - & 5017440542767728817653936668259736712851 & 9971295471872206320850346254728427595419 \\
5 & 16 & - & 8502629146414965332977658041614347092992 & 5079229625482867887148742152002214394652 \\
5 & 18 & - & 3618335194997566536226827750882872668071 & 1008167682823746101036115011159262928281 \\
5 & 20 & - & 6916057142617603142425932379858113204693 & 3353722048915161143290954452508953535871 \\
5 & 22 & - & 6229986141485489599799073235954532435191 & 5193506325962246267098254360183369410769 \\
5 & 24 & - & 8277133852911422387168980634459652665422 & 3626330407990258713029171125095287595440 \\
5 & 26 & - & 8956813423682488960953403152599740018168 & 1232511068285413761227747323600773966427 \\
5 & 28 & - & 6089044071216330152581858963214315608749 & 7892570466757351867226165509429765588799 \\
5 & 30 & - & 1270032909682992650957638704159867983648 & 4651047760256189098913151679708341583649 \\
5 & 32 & - & 1869892387750523118697231927844608165669 & 5753304515840789134956934603857679738268 \\
5 & 34 & - & 9684152719517059758944799921772022942830 & 7170072555964315759226822460244347476120 \\
5 & 36 & - & 1180843431364777370849783507791826395202 & 8228921484500748994185832814718042754140 \\
5 & 38 & - & 4836105411004967913248966188431293618832 & 9742989004710795938980186664367579667791 \\
5 & 40 & - & 0000000013323742450144361900466272570329 & 0000000000000000000000000000000000000000 \\
6 & 0 & - & 7394502272984717218669606554234844507336 & 7900747586357986312206316797987562404476 \\
6 & 2 & - & 0030626945829013820029345356214738247978 & 7097553551280981375209202414030805186526 \\
6 & 4 & - & 1406930893459767613755814019396522272257 & 5063292296670806011562154511959117376002 \\
6 & 6 & - & 9795378860969458764843473950612450501568 & 7753738011949868145765911242645683978731 \\
6 & 8 & - & 0762408088067160822192362017297286247520 & 1288019661542619671306015723197409273386 \\
6 & 10 & - & 6771011261025212236470592735664163594935 & 0310624815189162060223795361036128783042 \\
6 & 12 & - & 1413883010662720219367292601803594576682 & 2775941754043901346334229551487342801670 \\
6 & 14 & - & 7883257254414344931311965114352329096810 & 6312667269986962713706459780060161612670 \\
6 & 16 & - & 4044943673745282687464799481631488247796 & 4106559385352877036878156889694065531833 \\
6 & 18 & - & 0822257859259741579140961132225502461676 & 4213878956143176663888989862654352729766 \\
6 & 20 & - & 1511844765385135261603460339485034833124 & 5661037192232941652168490875004261142650 \\
6 & 22 & - & 7003910345142028190135505088263210232900 & 9890219137631379026147378535006695969922 \\
6 & 24 & - & 2158800702839114590817053638033458800789 & 5451059658232628971798014017978719608057 \\
6 & 26 & - & 3920561510067545268281883009401175901556 & 4271524721459013601966886576034003159897 \\
6 & 28 & - & 4323502196153391482466261031316264525078 & 8604108238399886125421593385438785031083 \\
6 & 30 & - & 5873628914233990714548803484464141395394 & 5537153623789338756155804010172708877892 \\
6 & 32 & - & 4554035429359391455890056711106780519789 & 3737564781384568742185602868292174749977 \\
6 & 34 & - & 5605725711896949918525042830316171528249 & 7978715927005890114618837470469463450989 \\
6 & 36 & - & 0125660383787731850306197219666274410303 & 2954327843721820795923850335185852046671 \\
6 & 38 & - & 7318757382688442061524835460951315603183 & 0593919366885833355745520122487143600971 \\
6 & 40 & - & 0000000000488697956771079805684279131680 & 0000000000000000000000000000000000000000 \\
7 & 0 & + & 8238382847645669449911056299829287728870 & 8118433851467157570429089515350680353155 \\
7 & 2 & + & 3589471250276311329938464089664563597105 & 4602957765579926763722264434781779691165 \\
7 & 4 & + & 1127057716992917539268710000204537862660 & 6231679317673998232343595990984583369547 \\
7 & 6 & + & 4319572554543241394449761863724641985470 & 7101258060926286535460971707904554670294 \\
7 & 8 & + & 5019983355599193844091072676223151284836 & 3255723154365949706700537372897069476759 \\
7 & 10 & + & 5281210383228254714964759271832000946490 & 8029015681518701833712836080378433743963 \\
7 & 12 & + & 1610355165570334192860025395211031687550 & 9787197874740303037683210168905769451538 \\
7 & 14 & + & 5484720932005410485519741764653009536832 & 7223182626981036772285702988619526017410 \\
7 & 16 & + & 6537244714163772913724947995315522429257 & 2231237330058874998929559492110579894829 \\
7 & 18 & + & 8422997159858024599898613134708304854075 & 5506963597438848868230793136335931976238 \\
7 & 20 & + & 0050992139281561627083575644501609604460 & 5450041428727690490127557492499187139131 \\
7 & 22 & + & 9820929343045999084888442480150952621522 & 5160277830001213114137968207757108681214 \\
7 & 24 & + & 7340823452749507454553410458038174354836 & 8642167322622639155762727208701017303322 \\
7 & 26 & + & 8943296813110844468009433966753845147584 & 1384290385791065761127279556016717886771 \\
7 & 28 & + & 2987879645220924186285642930101690587962 & 9393309315989535477502422859690150521686 \\
7 & 30 & + & 0616892429736085091217218193045074733730 & 6086804355652512678930525323770593822264 \\
7 & 32 & + & 4568873573158502457645987978839245045159 & 5089756837612652329928977614612053679054 \\
7 & 34 & + & 3688777553292144407500725597671282744002 & 0570352628378446119574051922953485737625 \\
7 & 36 & + & 3689272817044014581429930129543314879205 & 2274086482872868856684294558538594849570 \\
7 & 38 & + & 2334350978703060116515725051121653128202 & 6739680396271523088749712707922840990595 \\
7 & 40 & + & 0000000000065342587007460403399727124057 & 0000000000000000000000000000000000000000 \\
8 & 0 & + & 2354818596140470824446365031651667765457 & 0388411779551318823510289225055673749854 \\
8 & 2 & + & 3573652673484856427196329520860120015872 & 8150885013004075319656769768226301965581 \\
8 & 4 & + & 9574794651224645672316022742365441318801 & 1428758106680540155064067194103153970355 \\
8 & 6 & + & 1839649044588942778574477327937849440518 & 5015954064735708790521568195670769716134 \\
8 & 8 & + & 0228789264964763058741811018637036315871 & 8090890201107520116883320936938089255752 \\
8 & 10 & + & 3851982104340673708964967207774288336535 & 9000770806468739831379597891226394610394 \\
8 & 12 & + & 3474112789530240935899155179325232161728 & 2143748892824782018513079404734651773204 \\
8 & 14 & + & 9400284827050662524621669565101263132838 & 6721589071750406161046713738424491697078 \\
8 & 16 & + & 9171782049522307268389487713515414481377 & 5061038900391738249134870444612743024674 \\
8 & 18 & + & 6432361930740799118394970569832641340598 & 1465968667392373006767838585837187930512 \\
8 & 20 & + & 4308888509720205044867634044150985367044 & 9485980574608798257923560470179346044961 \\
8 & 22 & + & 1392638659486741007822312817494509014388 & 1509039128095550883156573189884882595120 \\
8 & 24 & + & 1423070101053334508566248465705318576161 & 4239617627930127841223304077074387107821 \\
8 & 26 & + & 5664578848348972703064126430460220142179 & 6502381898243167705061424808398326222061 \\
8 & 28 & + & 8048226764112735034339805841414111710072 & 6434966688650260596070223160921882387678 \\
8 & 30 & + & 1682141303651211353711803565327182092979 & 4172517607880170868071600440122325074703 \\
8 & 32 & + & 5555757165053622511155433529983438998094 & 5889590607172407550343108622431578370421 \\
8 & 34 & + & 3793657774466494635989336553352455342444 & 1761367584123254820455335571185935317344 \\
8 & 36 & + & 7505192168832973898801087123309253416478 & 5571329614083657956950133379156330171906 \\
8 & 38 & + & 4143376107418706949826706867003772521752 & 6816011113994672952308697237720661754903 \\
8 & 40 & + & 0000000000018068308301814838757301058353 & 0000000000000000000000000000000000000000 \\
9 & 0 & + & 2557518684441046362430552890298071220621 & 3866561037538909585815473902045196156104 \\
9 & 2 & + & 8375118361540559255050698073092512300039 & 8002195111136593176113542549193560266090 \\
9 & 4 & + & 0798510714269681079605852649875281358642 & 8937493275489261760753868232314117603636 \\
9 & 6 & + & 9372317674469764213049751162152022709991 & 7680916105847249122063175732111651819239 \\
9 & 8 & + & 7767945300733695831419676159600437634553 & 0247699516811377178849494109638111989367 \\
9 & 10 & + & 9955062939047907015093008018105664337851 & 8424370954581196253582418979121945549237 \\
9 & 12 & + & 4141160553101977931433015854621073476704 & 5815635785035407260332596081261138970277 \\
9 & 14 & + & 3049518455719809900705717129053306381864 & 2822332025930454713847187142184996041975 \\
9 & 16 & + & 1500041313840597953789571347209653659486 & 5930145846528578777986760475230088895254 \\
9 & 18 & + & 2400948176950921929370188402850395578506 & 1735459203928722026996267460318979246743 \\
9 & 20 & + & 8342210720914666485619798990030930812476 & 2188533155502557126847565228988457785658 \\
9 & 22 & + & 7414665215366333855683986312140792008281 & 7582541460785467184623289753704302385709 \\
9 & 24 & + & 8836713191487119389662094536143836612726 & 2152294420748825152486644928869356198920 \\
9 & 26 & + & 5620005061281153208019474314090086783777 & 3753089655774947884030807497391428674341 \\
9 & 28 & + & 5127239278323273444213430514405645025372 & 4924238436475234661798723207073075196206 \\
9 & 30 & + & 7291959158248927389366718701684607289397 & 7906293229033702002704580848313522755533 \\
9 & 32 & + & 3205221214029349777378046327585648231986 & 8727287311871342820826484825091914754652 \\
9 & 34 & + & 3680848775701684033806262188217204282994 & 9210897484390332893061566383738600827209 \\
9 & 36 & + & 1504853108886508221030200870082128104946 & 1704567539430775451002372990473279458945 \\
9 & 38 & + & 6044662247907894762126171100697383872845 & 2152937637825942225775428459576577452024 \\
9 & 40 & + & 0000000000002629076892843727929678227548 & 0000000000000000000000000000000000000000 \\
10 & 0 & + & 7353987436396797253628724489729302857755 & 5273403500536494234401536924078458584833 \\
10 & 2 & + & 1401127917321328559401196825666809273285 & 6315981283743436324465177189626097345261 \\
10 & 4 & + & 3373463248538246313301323971338151761118 & 8116342367682958197897456618304664106032 \\
10 & 6 & + & 3041397242717461841388684196623033622004 & 3117060556427604180741201634124361023969 \\
10 & 8 & + & 0943643532475361202445263202349734512849 & 3893493058215865506184535518774857221645 \\
10 & 10 & + & 8164053402431415136863060122068248514818 & 2074746967592066534796289044234098900700 \\
10 & 12 & + & 3477228668531847409916128433102738655961 & 4005431205998608532457473606630041724045 \\
10 & 14 & + & 2279001731902550657895749155504105447613 & 5213512655737020603344394806230564943953 \\
10 & 16 & + & 3523951808198425635099102195522511852857 & 1209494734674125524772827074198539749925 \\
10 & 18 & + & 1162723854529184955372870904211608437265 & 8278462054121871500184528414683394384746 \\
10 & 20 & + & 4918434325477973570079219531092301552544 & 2579110923203793306263902129105857511381 \\
10 & 22 & + & 8919744110506064538611222355037039930990 & 4724156936354691340925493126834525037416 \\
10 & 24 & + & 9223048441598762385906527220590988100734 & 9463101534220121666140877568111706105883 \\
10 & 26 & + & 9606292752276640879315851053471638578948 & 2121964904448883129904839737376110168492 \\
10 & 28 & + & 4484720556568965074515264922568477123950 & 3033773021174322221457729013140380077281 \\
10 & 30 & + & 5349719608582818935083265568666361149411 & 6496648317351270563291708346338463776301 \\
10 & 32 & + & 3254353987187937664725064966217369252154 & 7780706558722817226000333171435132837221 \\
10 & 34 & + & 2661925713335567593404972428903036059582 & 4030678956715165791627596733030254020225 \\
10 & 36 & + & 3757081892271528362198370608905198488995 & 3961880334644723103467983833088387469113 \\
10 & 38 & + & 3436156407606847826830776217587049351063 & 3430602205584376303999219106970022505905 \\
10 & 40 & + & 0000000000000281473394789457156348984914 & 0000000000000000000000000000000000000000 \\
11 & 0 & + & 6230259079649530845378347718381135133114 & 1075133513432887717382412164185262801654 \\
11 & 2 & + & 3925613306489854151531127857236258510516 & 6941429227603933407572450426411532098374 \\
11 & 4 & + & 6841275624459963720955079763943301147041 & 9084292622388980174940134210762883270229 \\
11 & 6 & + & 0681012873671684479347701609244146061591 & 0356333472936619888206063363367408716914 \\
11 & 8 & + & 2616161012511640407863860174350110684624 & 2068486347235793444918376196588113208983 \\
11 & 10 & + & 8486122757404765178286541753291046015206 & 2389682241397642053948925809955019585548 \\
11 & 12 & + & 5073089504710665612212469049079366869232 & 2394996185295722631413037677324109260678 \\
11 & 14 & + & 3340883612341516720388598194385511157433 & 4253086037750530869790473290155085357764 \\
11 & 16 & + & 0931900269296902838702249368524598268235 & 2723845899666256370156123522979283371627 \\
11 & 18 & + & 1360193058486939721451224488269977639754 & 2066698728271134089478552619115081902492 \\
11 & 20 & + & 8371663305021680462665594272712254599276 & 8472942361459686733450144494502578188482 \\
11 & 22 & + & 7498718802492795969015526345496593361522 & 0146221763807863934786199132932719278770 \\
11 & 24 & + & 2610672485808054601675289033817537212421 & 9678366616177831898724023351653932448639 \\
11 & 26 & + & 8709286258974696283667834956805982847399 & 5171089213032099789722202692180537400296 \\
11 & 28 & + & 8111022645304864742328886550991743706322 & 4825603082056196401039057739084247121624 \\
11 & 30 & + & 9895139856065386252208265777175683976965 & 3452528933612029445159595955968704694551 \\
11 & 32 & + & 6934539839814640270824033414728065432684 & 4566547151934018491821644846818371750436 \\
11 & 34 & + & 0826590848473693481242104820342034237560 & 6722964081431115379232494417569785243415 \\
11 & 36 & + & 3976835127961602937110709299632772715242 & 4428971908477581248544915294694676862005 \\
11 & 38 & + & 8134872566750411945116425534115768857661 & 9917085816514341966654131072720447865439 \\
11 & 40 & + & 0000000000000021932918118313448336161517 & 0000000000000000000000000000000000000000 \\
12 & 0 & + & 1191429103191373749682031604573321991180 & 8024765050558464430149045230609518073975 \\
12 & 2 & + & 0469064850413005852965094174787615016717 & 2974989178423205124834826959452437976170 \\
12 & 4 & + & 0605171427529337066754231435572396960708 & 0329134468021075124701427316048577041502 \\
12 & 6 & + & 0864982270853631048914727131772006359453 & 8760875082949084468594748953240027655389 \\
12 & 8 & + & 3401869394848613194684012295778245128572 & 6838943487074104866801721356989166558347 \\
12 & 10 & + & 0523106238080437677904097945797914557261 & 9126379239935411901912612928067913771880 \\
12 & 12 & + & 8964344185602225709120279820609767668694 & 9349933589792295204937002439177177904655 \\
12 & 14 & + & 6375798027542226632439672672550685682770 & 4716654713168773185428738431438233118701 \\
12 & 16 & + & 1038241022024003933757909135041487504320 & 1531245716588865592314894006234262479841 \\
12 & 18 & + & 4463986502048476346591660695510095572648 & 9738947914650978932904491079293813157390 \\
12 & 20 & + & 1022904479500434047264653514105248924927 & 7167395593295361621349662856919335261363 \\
12 & 22 & + & 6012127228891726906744793671674575037339 & 5983417023281281871567867461177589664551 \\
12 & 24 & + & 2711533354847296917122254923921673996009 & 3294707828310697681210072130016574014597 \\
12 & 26 & + & 6956013694527705478647737282919784220068 & 8696423104656137439421350436401949046774 \\
12 & 28 & + & 5337250643249176807447682919294346004090 & 7894333305832337497717626217184885048207 \\
12 & 30 & + & 1302290424020981971050610556795678855841 & 7283987634966711828697623927141837174110 \\
12 & 32 & + & 4623668671109060910756636236932390490455 & 6643996848072634325786362399676508277127 \\
12 & 34 & + & 2101433867944696686314753666248943334576 & 9819192864617659907436324967726417994845 \\
12 & 36 & + & 3217084489455528038452646868337685913777 & 6309469769274176669032291658466187619744 \\
12 & 38 & + & 3699669806495026242338123298551429984214 & 7249316270528151053055268737205697369448 \\
12 & 40 & + & 0000000000000000804471589762007412743324 & 0000000000000000000000000000000000000000 \\
13 & 0 & - & 4789234904214733368149926303181856782209 & 8907610975540148236942147166418679569385 \\
13 & 2 & - & 6037770799979869088401446720390934641623 & 0374805121640810263033757828116200302127 \\
13 & 4 & - & 3353677462293910167433327604661209034816 & 1099925358869439815682161299818056129614 \\
13 & 6 & - & 9817461157627510826240936483752487638277 & 6667998550070881386596611090964348583616 \\
13 & 8 & - & 7441949213334737948139775663762347045540 & 9453908432690903277354521474563830151918 \\
13 & 10 & - & 7400979058123439001563702106352001228823 & 5364730647138703856532766123624132857857 \\
13 & 12 & - & 6281902397180357847756085076601257059233 & 8826644341467547250884762227112243711582 \\
13 & 14 & - & 7395683363605230745847867733009691468229 & 1071195407445962630015440509302651341858 \\
13 & 16 & - & 2564790582378045159291922132801420701627 & 0999857785572890709765892995460163537746 \\
13 & 18 & - & 5873712820954344078615368785023873792967 & 1042428449990626175168597767981342664731 \\
13 & 20 & - & 7686371476201955959456736937674484685552 & 8429499826761278952287831611873696327533 \\
13 & 22 & - & 9491381007467180026389652703856318815841 & 7362655462707451746631557275235851930710 \\
13 & 24 & - & 7327020709093015659184397673723396489717 & 6557160426254224457588805301201132307814 \\
13 & 26 & - & 9619225532641447060602735676169375529145 & 6250937025642940845963050437259232434184 \\
13 & 28 & - & 5572243908573063271282309481379258382432 & 8708823562548138048816831598337804490207 \\
13 & 30 & - & 6957820716302490972982164623797766452344 & 1067757291009655258611624087748670379422 \\
13 & 32 & - & 8451282638444371647514061105914350471920 & 2511160955119262151364895881586509257520 \\
13 & 34 & - & 3663518281249890161994219173475294715484 & 3976158931822402826907269900543062867102 \\
13 & 36 & - & 4493077399338364321032325269091694264164 & 1826837897655945397183442571322125072356 \\
13 & 38 & - & 0839822179218487785899168252840206307686 & 4633723807611019631187336730426423472031 \\
13 & 40 & - & 0000000000000000107563893240661754162147 & 0000000000000000000000000000000000000000 \\
14 & 0 & - & 0452434226110034880994737455490328367891 & 4074109970497847999232979950374394705316 \\
14 & 2 & - & 6257130481022466721145338652509880685236 & 0036260346153016455520562630163147038335 \\
14 & 4 & - & 0024723244975079244471000161393013362892 & 4570549472668508878653752015232899323975 \\
14 & 6 & - & 0017030099095724920744639770443613877486 & 1502070748378166634795382180802370821352 \\
14 & 8 & - & 3597973085528712880453460805508527988126 & 1311707805801874540253519666496279904504 \\
14 & 10 & - & 3926385698415124856834892819505520076487 & 6533030721465264429739299537806277496792 \\
14 & 12 & - & 1932075938171188198836480717891397628618 & 8821734355944945002450421680946921053111 \\
14 & 14 & - & 0812170236452140960813636264123016951460 & 4090491622260801839636414599203428513399 \\
14 & 16 & - & 5768624077240560245880167843929878943362 & 6192590774128466958659967279906939091943 \\
14 & 18 & - & 4527727708914822737847794173707603983807 & 4719854577320232018867097298425688668578 \\
14 & 20 & - & 9559194972652041901898089536957783736061 & 7178461603345982165972489050106180224240 \\
14 & 22 & - & 6753038562725396529898715618066088612449 & 5486836957578979227454705312781939479776 \\
14 & 24 & - & 3919128867489080349250414435016185181570 & 5105104788686597661328731820644056539162 \\
14 & 26 & - & 0691706213613084656152673683781848556442 & 9083521751874921512979914667444529063696 \\
14 & 28 & - & 3257363264150749782919797508640854477047 & 0865430746155519067064137633853019137376 \\
14 & 30 & - & 9163715964320290640819249548948721004247 & 5229244986970967216842172633419540770894 \\
14 & 32 & - & 9501250120927964745212282372045789435653 & 4552445203132795427730753353342642061627 \\
14 & 34 & - & 0705916760904200708496527586737263382084 & 0323712749401392729048910828759008317001 \\
14 & 36 & - & 4224951315366490625825661223522389729953 & 9440801452524346536290631312108789095824 \\
14 & 38 & - & 5966040932792494378400525253938507189375 & 9083667907355352145253400435806902848376 \\
14 & 40 & - & 0000000000000000029743199255240346203575 & 0000000000000000000000000000000000000000 \\
15 & 0 & - & 8940193700433358460514604346420765417944 & 2062208532102458712759129313021455613646 \\
15 & 2 & - & 1852461950983703742024951621337080546294 & 7534814361731545073609092103657877558755 \\
15 & 4 & - & 6844043652639839190859289555427934406588 & 7053636005420111848083879696033684875196 \\
15 & 6 & - & 1120401501668010899979249937070815681081 & 9443337125228992231621634967382638091674 \\
15 & 8 & - & 3197127311594215974923585976336429557020 & 1830904861145164360926823073580636162028 \\
15 & 10 & - & 1562943487879333499830868301889075913563 & 5115001114889022406130372263331909533568 \\
15 & 12 & - & 3236747177668776479400277670079666805157 & 5099906270370014418968616101988624710522 \\
15 & 14 & - & 3121749995636093553324478941575688206221 & 2500206145161600316599437261937236858406 \\
15 & 16 & - & 3496149229860762616711316844015036379080 & 9718530956371109099569404876382135671131 \\
15 & 18 & - & 8042566382683273621307922213122858104685 & 1454399450303073806207113171813764966313 \\
15 & 20 & - & 2347785439814861721841398201047831263325 & 3253723032521865406062377047063752496268 \\
15 & 22 & - & 8362844638253951798932538568180118836215 & 4402711332047792729507645335745282096372 \\
15 & 24 & - & 7696460144504390062451348223401011857035 & 0558360011202231035207868870928175243596 \\
15 & 26 & - & 7144538712611161677362355621549534467178 & 5386190952982093136350388303700603716655 \\
15 & 28 & - & 6888260451134262557871440294672607078025 & 2670150645936795131466345211234946375305 \\
15 & 30 & - & 6924149259105256504185200135505573310911 & 0014985093690775221170171104122642017097 \\
15 & 32 & - & 7892863864637280782899040041637553296019 & 5568010687441832951601642985551937876483 \\
15 & 34 & - & 4041429036261240322002334015681315586517 & 2551454412263361202121996123007710385247 \\
15 & 36 & - & 1881790115289606886205277009565580343353 & 3025055948886788666030386470867041618943 \\
15 & 38 & - & 4247973491347983871490946986772158783461 & 8715837300407956058791361039834101290281 \\
15 & 40 & - & 0000000000000000004327862718246001981390 & 0000000000000000000000000000000000000000 \\
16 & 0 & - & 1109798837323746242968551453596445802206 & 4338366561630071686729586434523538844502 \\
16 & 2 & - & 9585762894230237319488590570752340547380 & 6566671317079902469554120416252805813491 \\
16 & 4 & - & 4515190360995157289278440970672763391022 & 9167721075956422115632268345557484956464 \\
16 & 6 & - & 5608860029031504518552988845902280444503 & 9834458251195783249646111827783911421717 \\
16 & 8 & - & 6476637150588070942865794835291984091460 & 0789562774659848101644596221536941537853 \\
16 & 10 & - & 6854954072538063625712239589584976924210 & 1713353217644655995928406931263799572283 \\
16 & 12 & - & 4320375570511299256983954303616078931748 & 2499503181677368923195775480976629436799 \\
16 & 14 & - & 7415516496276409505073703118401962967204 & 1482617810256008512923791257383317983033 \\
16 & 16 & - & 4989666812673890149655277951818280546098 & 0766947108120717526618903042747957332922 \\
16 & 18 & - & 9926394015910401002470583904663669350827 & 9653758433336746095142896211070593686243 \\
16 & 20 & - & 2821560736701829232666203498054148669673 & 8194722823588827606683601735420460926235 \\
16 & 22 & - & 0230700947334815135237758024420930054109 & 6713456216194874174036769882285129279422 \\
16 & 24 & - & 2079478269966985003967667527809452413373 & 3471547828129228783730218088680119254478 \\
16 & 26 & - & 8637715525941685364578178299790358891780 & 8093763821340022536279839174353865899946 \\
16 & 28 & - & 6500216546310658574430197907963650852394 & 1119819487668177013818853616048795989205 \\
16 & 30 & - & 4894852108828465801061747702089300636906 & 1444374636046332420287050973553094374923 \\
16 & 32 & - & 4087978055296074656045042806872435086691 & 8320850337246429776628694902369329059063 \\
16 & 34 & - & 1642225568187657205075507655221535768841 & 9866138884421289908363704595594966040455 \\
16 & 36 & - & 2587168531962766630958481732493426748290 & 4752569286990688178083050598953035275326 \\
16 & 38 & - & 8968862005559874188760346974114921569979 & 4391631078960698045299293606808446488320 \\
16 & 40 & - & 0000000000000000000463348262960005594473 & 0000000000000000000000000000000000000000 \\
17 & 0 & - & 6361250453941912458867122414602821534195 & 2861244228685644896442731557023963313420 \\
17 & 2 & - & 5235477952572537597819139701326181288906 & 0978677101815753039029275239404149413426 \\
17 & 4 & - & 1468816394023268056415363464201058885542 & 3934520420427372329293901627292424322415 \\
17 & 6 & - & 1060831318888252948333009365290167484802 & 7504184543780842489309034373431371840946 \\
17 & 8 & - & 4810036206415250953773749966689013165618 & 1592920057730474581012726645612324836366 \\
17 & 10 & - & 2539617471761158507944984561124468007405 & 3863962166987699501460049280483811955079 \\
17 & 12 & - & 5242026763913335611495000766843702354057 & 9543545634008682667761603563996918852388 \\
17 & 14 & - & 0236588433629775311027776163465142721006 & 0046227865090272555053985974940165779793 \\
17 & 16 & - & 7089614959973531532558325587048729191383 & 8115632610509134155678368736346733231112 \\
17 & 18 & - & 3307491655215782707881640404532113183186 & 5936357050459752541672192923244997218894 \\
17 & 20 & - & 7848499792819211779938952037950461624124 & 5047590875560328253031764088744032236484 \\
17 & 22 & - & 1173025864233250866354964490945439264506 & 3977131534574329482974732706639342123725 \\
17 & 24 & - & 8040430088564720765307295731553555017512 & 8226384895836777329362268898407334593174 \\
17 & 26 & - & 9159991636620500456926760594943341609055 & 5814784322554446264895377754534498941924 \\
17 & 28 & - & 0320600693539272498495979158283602981111 & 8287293477683740261132203006380991283719 \\
17 & 30 & - & 8326073109017999571578416501688372422218 & 6912873587295394081380856319161856510675 \\
17 & 32 & - & 9830323765239680801974267451572776348798 & 0068524121908800925834608648406211904213 \\
17 & 34 & - & 8421051264139757967374339165957804271579 & 2025841956661308085366256004498664149366 \\
17 & 36 & - & 4958095250747446761018760374161331217430 & 6206317576439170368344401862742528337221 \\
17 & 38 & - & 1192983921029765518598352558328645943410 & 5523721663722388372926529015510334613895 \\
17 & 40 & - & 0000000000000000000036104938157178965720 & 0000000000000000000000000000000000000000 \\
\end{longtable}
\endgroup

